\documentclass[12pt,a4paper]{article}
\usepackage[UTF8]{ctex}
\usepackage{amsmath,amssymb,amsfonts}
\usepackage{graphicx}
\usepackage{booktabs}
\usepackage[margin=2.5cm]{geometry}

\title{多时相遥感影像变化检测算法说明}
\author{}
\date{}

\begin{document}
\maketitle

\section{RMFS 黎曼流形特征空间分类}

\subsection{概述}

RMFS（Riemannian Manifold Feature Space）模块在 CNN 主干网络提取的特征图基础上，
通过协方差池化将特征映射到对称正定（SPD）流形空间，
再利用 Log-Euclidean 映射将流形特征投射到切空间进行欧氏分类。
该方法能够捕获特征通道间的二阶统计信息，相比传统全局平均池化具有更强的特征表达能力。

\subsection{处理流程}

完整的 RMFS 处理流程如下：

\begin{enumerate}
    \item \textbf{CNN 提取局部特征}：主干网络（VGG/ResNet/DenseNet/MobileNet/EfficientNet）
          提取输入图像的特征图 $\mathbf{X} \in \mathbb{R}^{B \times C \times H \times W}$。
    \item \textbf{协方差池化}：对特征图进行通道归一化和去中心化后，
          计算协方差矩阵 $\mathbf{\Sigma} = \frac{1}{N-1}(\mathbf{X} - \bar{\mathbf{X}})(\mathbf{X} - \bar{\mathbf{X}})^\top$，
          其中 $N = H \times W$，得到 $\mathbf{\Sigma} \in \mathbb{R}^{B \times C \times C}$。
    \item \textbf{SPD 流形表示}：对协方差矩阵加入正则化项 $\mathbf{\Sigma} \leftarrow \mathbf{\Sigma} + \epsilon \mathbf{I}$，
          确保矩阵严格正定，位于 SPD 流形 $\mathcal{S}_{++}^C$ 上。
    \item \textbf{Log-Euclidean 映射}：通过特征分解 $\mathbf{\Sigma} = \mathbf{V} \text{diag}(\lambda_i) \mathbf{V}^\top$，
          计算矩阵对数 $\log(\mathbf{\Sigma}) = \mathbf{V} \text{diag}(\log \lambda_i) \mathbf{V}^\top$，
          将 SPD 矩阵映射到切空间。
    \item \textbf{欧氏分类器}：将对数映射后的矩阵展平为向量 $\mathbf{v} \in \mathbb{R}^{C^2}$，
          经 Dropout 正则化后送入全连接层进行分类。
\end{enumerate}

\subsection{通道降维}

当主干网络输出通道数 $C > 512$ 时（如 ResNet50 的 2048、MobileNet 的 1280），
RMFS 模块会先通过 $1 \times 1$ 卷积将通道数降至 256，
以避免协方差矩阵维度过大导致的计算和内存开销。

\subsection{兼容性}

RMFS 模块通过 \texttt{use\_rmfs} 参数控制开关：
\begin{itemize}
    \item \texttt{use\_rmfs=True}：CNN $\rightarrow$ RMFS $\rightarrow$ Linear 分类器
    \item \texttt{use\_rmfs=False}：CNN $\rightarrow$ GAP $\rightarrow$ Linear 分类器（原始流程）
\end{itemize}

\end{document}
